% GENERATED FILE. Edit canonical lessons, then run npm run generate:books.
% canonical-source-hash: fnv1a64:a2e58f84520fad9a
% canonical-lessons: SA-C197-venuvanam, SA-C197-indradhanuh, SA-C197-karaka, SA-C197-avasyayah, SA-C197-jyotsna

\chapter{Bamboo grove, Rainbow, Hail, Dew, Moonlight}
\label{ch:sa-bamboo-grove-rainbow-hail-dew-moonlight}
\hlchaptermodality{197}

\begin{chapteropening}
\textbf{By the end of this chapter:} \emph{I can say a bamboo grove, a rainbow, hail, dew and moonlight.}

Complete the last lesson of chapter 197: I can say a bamboo grove, a rainbow, hail, dew and moonlight.
\end{chapteropening}

\section[\texorpdfstring{\sk{वेणुवनम्} (veṇuvanam)}{वेणुवनम् (veṇuvanam)}]{\sk{वेणुवनम्} (veṇuvanam) — a bamboo grove}

\label{lesson:SA-C197-venuvanam}



\begin{quote}
Before the new one: say the Sanskrit for a stream, a river, then the Sanskrit for a cave.
\end{quote}

\subsection*{You'll want to know: \sk{वेणुवनम्}}
\textbf{\sk{वेणुवनम्}} — \emph{veṇuvanam} — \textquotedblleft{}a bamboo grove\textquotedblright{}.

\begin{grammarlens}[title={What you've built}]
One more word for this chapter.
\end{grammarlens}

\subsection*{Your turn}
\begin{itemize}
\raggedright
  \item \emph{Say it:} \emph{veṇuvanam}
  \item \emph{Say it:} \emph{veṇuvanam}, once more
  \item \emph{Say it:} say \emph{veṇuvanam}
  \item \emph{Recall:} say the Sanskrit for a stream, a river, then the Sanskrit for a cave, then say \emph{veṇuvanam} again
\end{itemize}

\begin{tcolorbox}[breakable,colback=teal!4,colframe=teal!35!black,title={Before you move on}]
What does \sk{वेणुवनम्} mean? (\textquotedblleft{}A bamboo grove\textquotedblright{}.) Say it once more.
\end{tcolorbox}
\section[\texorpdfstring{\sk{इन्द्रधनुः} (indradhanuḥ)}{इन्द्रधनुः (indradhanuḥ)}]{\sk{इन्द्रधनुः} (indradhanuḥ) — a rainbow}

\label{lesson:SA-C197-indradhanuh}



\begin{quote}
Before the new one: say the Sanskrit for a cave, then the Sanskrit for a bamboo grove.
\end{quote}

\subsection*{You'll want to know: \sk{इन्द्रधनुः}}
\textbf{\sk{इन्द्रधनुः}} — \emph{indradhanuḥ} — \textquotedblleft{}a rainbow\textquotedblright{}.

\begin{grammarlens}[title={What you've built}]
One more word for this chapter.
\end{grammarlens}

\subsection*{Your turn}
\begin{itemize}
\raggedright
  \item \emph{Say it:} \emph{indradhanuḥ}
  \item \emph{Say it:} \emph{indradhanuḥ}, once more
  \item \emph{Say it:} say \emph{indradhanuḥ}
  \item \emph{Recall:} say the Sanskrit for a cave, then the Sanskrit for a bamboo grove, then say \emph{indradhanuḥ} again
\end{itemize}

\begin{tcolorbox}[breakable,colback=teal!4,colframe=teal!35!black,title={Before you move on}]
What does \sk{इन्द्रधनुः} mean? (\textquotedblleft{}A rainbow\textquotedblright{}.) Say it once more.
\end{tcolorbox}
\section[\texorpdfstring{\sk{करका} (karakā)}{करका (karakā)}]{\sk{करका} (karakā) — hail}

\label{lesson:SA-C197-karaka}



\begin{quote}
Before the new one: say the Sanskrit for a bamboo grove, then the Sanskrit for a rainbow.
\end{quote}

\subsection*{You'll want to know: \sk{करका}}
\textbf{\sk{करका}} — \emph{karakā} — \textquotedblleft{}hail\textquotedblright{}.

\begin{grammarlens}[title={What you've built}]
One more word for this chapter.
\end{grammarlens}

\subsection*{Your turn}
\begin{itemize}
\raggedright
  \item \emph{Say it:} \emph{karakā}
  \item \emph{Say it:} \emph{karakā}, once more
  \item \emph{Say it:} say \emph{karakā}
  \item \emph{Recall:} say the Sanskrit for a bamboo grove, then the Sanskrit for a rainbow, then say \emph{karakā} again
\end{itemize}

\begin{tcolorbox}[breakable,colback=teal!4,colframe=teal!35!black,title={Before you move on}]
What does \sk{करका} mean? (\textquotedblleft{}Hail\textquotedblright{}.) Say it once more.
\end{tcolorbox}
\section[\texorpdfstring{\sk{अवश्यायः} (avaśyāyaḥ)}{अवश्यायः (avaśyāyaḥ)}]{\sk{अवश्यायः} (avaśyāyaḥ) — dew}

\label{lesson:SA-C197-avasyayah}



\begin{quote}
Before the new one: say the Sanskrit for a rainbow, then the Sanskrit for hail.
\end{quote}

\subsection*{You'll want to know: \sk{अवश्यायः}}
\textbf{\sk{अवश्यायः}} — \emph{avaśyāyaḥ} — \textquotedblleft{}dew\textquotedblright{}.

\begin{grammarlens}[title={What you've built}]
One more word for this chapter.
\end{grammarlens}

\subsection*{Your turn}
\begin{itemize}
\raggedright
  \item \emph{Say it:} \emph{avaśyāyaḥ}
  \item \emph{Say it:} \emph{avaśyāyaḥ}, once more
  \item \emph{Say it:} say \emph{avaśyāyaḥ}
  \item \emph{Recall:} say the Sanskrit for a rainbow, then the Sanskrit for hail, then say \emph{avaśyāyaḥ} again
\end{itemize}

\begin{tcolorbox}[breakable,colback=teal!4,colframe=teal!35!black,title={Before you move on}]
What does \sk{अवश्यायः} mean? (\textquotedblleft{}Dew\textquotedblright{}.) Say it once more.
\end{tcolorbox}
\section[\texorpdfstring{\sk{ज्योत्स्ना} (jyotsnā)}{ज्योत्स्ना (jyotsnā)}]{\sk{ज्योत्स्ना} (jyotsnā) — moonlight}

\label{lesson:SA-C197-jyotsna}



\begin{quote}
Before the new one: say the Sanskrit for hail, then the Sanskrit for dew.
\end{quote}

\subsection*{You'll want to know: \sk{ज्योत्स्ना}}
\textbf{\sk{ज्योत्स्ना}} — \emph{jyotsnā} — \textquotedblleft{}moonlight\textquotedblright{}.

\begin{grammarlens}[title={What you've built}]
One more word for this chapter.
\end{grammarlens}

\subsection*{Your turn}
\begin{itemize}
\raggedright
  \item \emph{Say it:} \emph{jyotsnā}
  \item \emph{Say it:} \emph{jyotsnā}, once more
  \item \emph{Say it:} say \emph{jyotsnā}
  \item \emph{Recall:} say the Sanskrit for hail, then the Sanskrit for dew, then say \emph{jyotsnā} again
\end{itemize}

\begin{tcolorbox}[breakable,colback=teal!4,colframe=teal!35!black,title={Before you move on}]
What does \sk{ज्योत्स्ना} mean? (\textquotedblleft{}Moonlight\textquotedblright{}.) Say it once more.
\end{tcolorbox}
